\section{Quality-guided maintenance and its assumptions}\label{sec:C15:maintenance}
An accumulating skill bank can become redundant, stale, and internally inconsistent. HyperSkill periodically prunes low-utility knowledge and merges structurally related skills~\cite{xu2026hyperskill}. Its inspected code computes a utility score combining observed success rate with normalized step efficiency. For retrieval count $\nu$, success count $\sigma$, mean steps $\bar T$, and weight $\beta$, the form is
\begin{equation}
 \gamma=\beta\frac{\sigma}{\nu}
 +(1-\beta)\left(1-\frac{\bar T-T_{\min}}{T_{\max}-T_{\min}}\right),
\end{equation}
for $\nu>0$, $\beta\in[0,1]$, $0<\bar T$, and $T_{\min}\leq\bar T\leq T_{\max}$ with $T_{\max}>T_{\min}$, where the last two quantities delimit the step-count normalization range. The implementation handles untested nodes and degenerate ranges separately~\cite{repo_hyperskill2026}. This score is a maintenance heuristic. It is not an unbiased estimate of a skill's individual treatment effect.

The attribution problem is concrete. When several skills are retrieved together and a task succeeds, updating every retrieved node credits them jointly. A harmful item can survive because other guidance compensates for it. A useful item can be penalized on intrinsically difficult tasks. The quality estimate is policy- and task-distribution-dependent. If the model improves or the task order changes, old scores need not remain calibrated. Maintenance should therefore retain exposure counts and evaluation conditions rather than only a scalar utility.

The inspected \path{update_utility} updates retrieved nodes and their incident hyperedges after an episode. Its comment specifies that the utility is consumed by maintenance. The ranking routine still uses co-occurrence and similarity, not utility as its primary ordering~\cite{repo_hyperskill2026}. This separation is easy to miss in a high-level summary. Changing the utility formula affects which knowledge survives and merges; it does not simply replace the online ranker with a quality score.

The normalization has a time-dependent implementation boundary. The core maintains global historical extrema of observed mean step counts, but stores each node's utility when that node is updated; pruning later reads that stored value. Without a full rescore, nodes can therefore carry utilities computed under different ranges. This is a source-level observation, not a demonstrated benchmark error. A discriminating test holds one node's exposure and outcome counts fixed, changes the global range through an unrelated trajectory, and compares cached with recomputed utility before pruning. Versioning the scoring context or recomputing all eligible scores are possible interventions.

Structure-informed merging uses quality-weighted co-occurrence and propagated embeddings to propose related abstractions. A merged node can reduce redundancy, but it can also erase incompatible preconditions. The merge model should not decide compatibility solely from a high cosine score. Two procedures can use similar vocabulary while applying to different authorization scopes or tool versions. Retain source nodes or reversible merge metadata when the cost of a bad abstraction is high.

Maintenance frequency creates another tradeoff. Frequent updates can respond quickly to drift but add embedding, graph, and language-model work. Infrequent updates let redundant or harmful items influence more tasks. A maintenance schedule chosen on one benchmark order can overfit that order. Evaluate shuffled sequences, abrupt environment changes, and mixtures of easy and difficult tasks while charging all extraction and merge costs. Without these controls, a smaller final graph can look efficient even if its construction consumed substantial resources.

The practical lesson is to treat consolidation as a state transition with evidence and rollback. Pruned knowledge may still be relevant to historical audit, while merged knowledge should preserve references to the experiences that justified it. A physically compact bank is not necessarily a semantically safer bank.

Complexity depends on implementation rather than the hypergraph label. The inspected core constructs a dense quality-weighted co-occurrence matrix for the eligible abstraction nodes~\cite{repo_hyperskill2026}. For $n$ such nodes, matrix storage alone is $O(n^2)$ before propagated embeddings and merge candidates. If each of $m$ episodes contains at most $r$ eligible nodes, forming all within-episode pairs costs up to $O(mr^2)$ under a direct construction. These are analytic upper bounds for the stated representation, not timings. Sparse incidence operations can change the practical tradeoff, but a sparse source graph does not guarantee that the derived co-occurrence matrix stays sparse.

This matters when a memory method claims online scalability. Retrieval can remain fast while periodic maintenance becomes the dominant cost. The maintenance interval amortizes that work across tasks but also delays removal of stale knowledge. A benchmark ending just before maintenance can undercount the eventual cost; one ending just after it can overstate per-query latency if the work is actually asynchronous. Report cumulative maintenance time and resource use over the same task horizon, alongside the critical-path delay experienced by each task.

